\subsection{问题三：实测条件推荐与缺测核查}

\subsubsection{收率与实测候选域}

组合 \(i\) 在温度 \(T\) 下的 C4 烯烃收率仍按式\eqref{eq:c4-yield}计算。第三问不拟合连续最优点，
而是在以下两个实测候选域内分别比较收率。

一般情景的实测候选域 \(D_{0}^{\mathrm{all}}\) 包含附件1全部114个组合—温度单元；
严格低温实测候选域为

\begin{equation}
D_{0}^{\mathrm{low}}
=\{(i,T)\in D_{0}^{\mathrm{all}}:T<350~^\circ\mathrm{C}\},
\tag{19}
\label{eq:q3-domain}
\end{equation}

共73个单元。一般情景取 \(D_0=D_{0}^{\mathrm{all}}\)，严格低温情景取
\(D_0=D_{0}^{\mathrm{low}}\)，两种情景分别求

\begin{equation}
(i^*,T^*)=\mathop{\arg\max}_{(i,T)\in D_0}Y_i(T).
\tag{20}
\label{eq:q3-observed-argmax}
\end{equation}

式\eqref{eq:q3-observed-argmax}给出的是附件实测范围内的最高观测收率，不是连续温度
上的全局最优，也不表示未来重复实验中必然仍为第一。

\subsubsection{两种温度情景的推荐结果}

表\ref{tab:q3-recommendations}列出两个实测候选域内的最高收率条件及其与次高值的差距。
\begin{table}[H]
\centering
\caption{两种情景下的实测条件推荐与主要差距}
\label{tab:q3-recommendations}
\begin{tabular}{lrrrr}
\toprule
情景 & 推荐组合 & 温度（℃） & 收率（\%） & 与次高差距（百分点） \\
\midrule
一般情景 & A3 & 400 & 44.7281 & 1.6096 \\
严格低于350~℃ & A2 & 325 & 17.2643 & 6.4716 \\
\bottomrule
\end{tabular}
\end{table}

一般情景的次高单元为A3在450~℃，收率43.1184\%；最高的不同组合为A4在400~℃，
收率36.2778\%，比推荐低8.4502个百分点。因此，一般情景的前两个观测都来自A3，
但A3在450~℃只是单组合个案，不能据此归纳所有组合的高温变化。

严格低温情景的次高单元也是最高的不同组合，即A3在325~℃，收率10.7927\%，比A2
低6.4716个百分点。由于每个条件没有重复实验，这些差距是观测差，而不是统计显著性
结论。

\subsubsection{内部缺测条件的辅助核查}

附件出现的温度标签为250、275、300、325、350、400和450~℃。只有A6--A14、B1和
B2缺少的325~℃同时位于各组合自身实测温度范围内部，共11个单元。本文对每个组合
分别用相邻已测点构造分段三次曲线，并限制曲线不产生与局部数据形态不符的额外振荡；
这种方法称为保形分段三次插值。计算不跨组合借用曲线，不外推，也不搜索标签之间的
任意连续温度。

验证时依次暂时移除一个已知的内部观测，用其余观测预测被移除点，再比较预测值与
真实值；这一操作称为内部遮点。该插值在72次内部遮点中的平均绝对误差为0.7152个
百分点，中位绝对误差为0.1627个
百分点，均低于局部直线的1.5570和0.8092个百分点；72次中有68次误差更小。但它在
5个温度块中只正确识别4次最高候选，最坏遮点误差仍为13.4906个百分点。因此，这一
方法的平均预测偏差和典型预测偏差虽然较小，个别点仍可能出现较大误差。故该方法只
用于缺测核查，预测值不作为观测或正式推荐。

11个缺口中，分段三次插值的最高预测为A7在325~℃的6.5517\%；改用局部直线后，
最高预测为7.2293\%。两者都低于严格低温实测推荐的17.2643\%，故低温推荐不依赖
这两种插值形式中的某一种。完整预测表保留在复现材料中，正文不逐项展开。

\subsubsection{问题三结论与适用范围}

325~℃只有10种组合具有实测值，450~℃只有A3具有实测值，说明附件在高温个案和
低温横向覆盖上都存在限制。遮点误差是预测验证误差，不是实验测量误差；现有无重复
数据也不支持传统显著性检验、置信区间或因果判断。

综上，一般情景推荐A3在400~℃的实测条件，严格低温情景推荐A2在325~℃的实测条件。
对11个已有标签缺口的辅助核查没有发现会改变低温推荐的候选。结论只适用于附件中的
组合和实测温度标签，不能外推到连续温度最优、新配方、工业规模、长期稳定性、安全性
或经济性。
